\documentclass[11pt,a4paper]{article}
\usepackage[T1]{fontenc}
\usepackage[utf8]{inputenc}
\usepackage[french]{babel}
\usepackage{lmodern,amsmath,amssymb,booktabs}
\usepackage[margin=25mm]{geometry}
\usepackage{microtype}
\usepackage[hidelinks]{hyperref}
\newcommand{\jj}{\mathrm{j}}
\newcommand{\Reel}{\operatorname{Re}}
\title{De la tension $v(t)$ au coefficient $S_{11}$\\Démodulation, séparation directionnelle et calibration}
\author{Projet MicroWaveTomography --- banc plexiglas}
\date{18 septembre 2026}
\begin{document}
\maketitle
\begin{abstract}
Nous détaillons ce que la démodulation en phase et en quadrature permet de
calculer, puis l'information supplémentaire nécessaire pour déterminer
le coefficient de réflexion. Une tension totale unique ne suffit pas à
séparer deux ondes de même fréquence. Nous distinguons donc le modèle
physique idéal, les mesures complexes fournies par l'instrument et leur
correction par calibration SOL. Les calculs ne supposent pas que le logiciel
Python reçoit des échantillons de la sinusoïde radiofréquence.
\end{abstract}

\section{Conventions : une fréquence fixée \texorpdfstring{$f_0$}{f0}}
On se place en régime sinusoïdal établi, après le transitoire de mise en
fréquence. L'enveloppe complexe $\alpha(t)$ est définie par rapport à la
porteuse de fréquence $f_0$, selon la convention
\begin{equation}
 v(t)=\Reel\{\alpha(t)\,e^{\jj2\pi f_0t}\},
 \qquad \alpha(t)=Ae^{\jj\phi}.
\end{equation}
Ainsi $v(t)=A\cos(2\pi f_0t+\phi)$ ; $A$ est une amplitude crête.
L'enveloppe complexe $\alpha(t)$ est constante pendant cette mesure
idéale : la notation $\alpha(t)$ désigne l'enveloppe en général, mais
elle ne dépend pas de $t$ dans le régime sinusoïdal établi considéré ici.
La lettre $t$ désigne ici le temps physique d'observation de la
sinusoïde, et non un retard de propagation à estimer.

Pour deux contributions de même fréquence,
\begin{equation}
 v(t)=A_1\cos(2\pi f_0t+\phi_1)+A_2\cos(2\pi f_0t+\phi_2),
 \qquad
 \alpha(t)=A_1e^{\jj\phi_1}+A_2e^{\jj\phi_2}.
\end{equation}
La démodulation fournit l'enveloppe complexe du signal par rapport à la
fréquence $f_0$, après mélange en phase et en quadrature, puis filtrage
passe-bas (voir section~\ref{s:demod}). Elle ne fournit pas
séparément ses deux termes. Les deux contributions ne sont pas séparables
par un filtrage en fréquence puisqu'elles ont la même fréquence.

\section{Calcul de \texorpdfstring{$I$ et $Q$ à partir de $v(t)$}{I et Q à partir de v(t)}}
\label{s:demod}
\subsection{Multiplication et filtrage passe-bas}
En écrivant $\alpha(t)=I+\jj Q$, on obtient
\begin{equation}
 v(t)=I\cos(2\pi f_0t)-Q\sin(2\pi f_0t).
\end{equation}
On multiplie par deux références cohérentes avec la source :
\begin{align}
 2v(t)\cos(2\pi f_0t)
 &=I+I\cos(4\pi f_0t)-Q\sin(4\pi f_0t),\\
 -2v(t)\sin(2\pi f_0t)
 &=Q-Q\cos(4\pi f_0t)-I\sin(4\pi f_0t).
\end{align}
Un filtre passe-bas élimine les termes à $2f_0$ et conserve les constantes :
\begin{equation}
 \boxed{I=\mathrm{PB}\{2v(t)\cos(2\pi f_0t)\},\qquad
 Q=\mathrm{PB}\{-2v(t)\sin(2\pi f_0t)\}.}
\end{equation}
Le signe moins dans la voie quadrature résulte de notre convention
$e^{+\jj2\pi f_0t}$. Une autre convention est possible si elle est employée
partout de manière cohérente.

\subsection{Expression par intégration}
Pour une durée $T$ contenant un nombre entier de périodes,
\begin{align}
 I&=\frac{2}{T}\int_{t_0}^{t_0+T}v(t)\cos(2\pi f_0t)\,\mathrm dt,\\
 Q&=-\frac{2}{T}\int_{t_0}^{t_0+T}v(t)\sin(2\pi f_0t)\,\mathrm dt,\\
 \alpha(t)=I+\jj Q
 &=\frac{2}{T}\int_{t_0}^{t_0+T}v(t)e^{-\jj2\pi f_0t}\,\mathrm dt.
\end{align}
Ces égalités sont exactes pour les sinusoïdes stationnaires considérées.
Sinon, le rejet du terme à $2f_0$ dépend du filtre ou de la fenêtre
 d'intégration. Une erreur de fréquence entre signal et référence produit
une rotation de l'enveloppe : $I$ et $Q$ ne sont alors plus constants.
Le bruit entraîne également des fluctuations autour de leurs valeurs moyennes.

\subsection{Expression discrète et lien avec la DFT}
Avec $N$ échantillons réels $v_n=v(t_n)$, $t_n=n/f_s$, une estimation est
\begin{equation}
 \widehat{\alpha}(t)=\frac{2}{N}\sum_{n=0}^{N-1}
 v_n e^{-\jj2\pi f_0 n/f_s}.
\end{equation}
Si $f_0=\ell_0f_s/N$, c'est $2/N$ fois le coefficient de DFT d'indice
$\ell_0$. Pour retrouver exactement l'amplitude complexe d'une sinusoïde
réelle, on suppose un échantillonnage cohérent, sans repliement, et une
fréquence distincte de zéro et de Nyquist.

Il s'agit d'une \emph{représentation mathématique de la démodulation}.
Un VNA peut d'abord transposer le signal vers une fréquence intermédiaire,
puis le numériser. Nous n'affirmons pas que le LiteVNA échantillonne
directement la radiofréquence ni qu'il utilise exactement cette somme.

\section{Pourquoi une tension totale ne suffit pas}
À un plan de référence donné sur une ligne d'impédance réelle
$Z_0=50\,\Omega$, notons $V^+$ et $V^-$ les amplitudes complexes des ondes
incidente et réfléchie. La tension totale y est
\begin{equation}
 v(t)=\Reel\{\alpha(t)e^{\jj2\pi f_0t}\},
 \qquad \alpha(t)=V^++V^-.
\end{equation}
La mesure de $\alpha(t)$ fournit deux nombres réels ; les deux amplitudes
complexes inconnues en représentent quatre. Pour tout nombre complexe $C$,
le couple $(V^++C,V^--C)$ donne la même tension totale.

\subsection{Si l'onde incidente est connue}
Si $V^+$ est connu \textbf{en amplitude et en phase, au même plan de
référence}, avec $V^+\ne0$, alors
\begin{equation}
 V^-=\alpha(t)-V^+,
 \qquad
 \boxed{\Gamma=\frac{V^-}{V^+}=\frac{I+\jj Q}{V^+}-1.}
\end{equation}
Connaître uniquement la puissance réglée sur le générateur ne fournit pas
cette amplitude complexe au plan de mesure.

\subsection{Si la tension démodulée est déjà celle de l'onde réfléchie}
Si un dispositif directionnel fournit un signal proportionnel à $V^-$,
et si ce facteur de proportionnalité est corrigé, alors la démodulation
fournit $I_-+\jj Q_-=V^-$. Dans ce cas,
\begin{equation}
 \boxed{\Gamma=\frac{I_-+\jj Q_-}{V^+}.}
\end{equation}
Il n'y a pas de terme $-1$ ici. \textbf{Il faut donc toujours préciser ce
que représente la tension $v(t)$ : tension totale ou signal de réflexion séparé.}

\section{L'information directionnelle supplémentaire}
Notons $\widetilde J$ l'amplitude complexe du courant, compté positivement
vers la charge. Nous utilisons $J$ pour éviter de confondre le courant avec
la composante en phase $I$. Pour les deux ondes,
\begin{equation}
 \alpha(t)=V^++V^-,\qquad
 Z_0\widetilde J=V^+-V^-.
\end{equation}
Le signe moins du courant réfléchi traduit son sens de propagation opposé.
Les deux équations donnent
\begin{equation}
 \boxed{V^+=\frac{\alpha(t)+Z_0\widetilde J}{2},\qquad
 V^-=\frac{\alpha(t)-Z_0\widetilde J}{2}.}
\end{equation}
Par conséquent,
\begin{equation}
 \boxed{\Gamma=
 \frac{\alpha(t)-Z_0\widetilde J}
      {\alpha(t)+Z_0\widetilde J}.}
\end{equation}
Un pont directionnel ou un coupleur directionnel permet d'obtenir des
combinaisons électriques sensibles au sens de propagation. Les équations
ci-dessus expliquent le principe de séparation ; elles ne signifient pas
que le LiteVNA mesure directement et séparément la tension et le courant
avec deux capteurs idéaux. Son schéma interne détaillé n'est pas établi ici.

Le manuel du LiteVNA indique l'emploi d'un mélangeur unique et de
commutations RF internes. Deux informations complexes ne nécessitent donc
pas nécessairement deux chaînes de démodulation simultanées. En cas de
mesure successive, leur comparaison suppose une stabilité et une cohérence
suffisantes pendant la séquence.

\subsection{Pourquoi le sens de propagation change le signe du courant}
La tension seule donne une somme : une information supplémentaire est
nécessaire. Prenons un axe $x$ dirigé vers l'antenne et un courant positif
dans ce même sens. Sur une ligne idéale sans pertes, les deux ondes peuvent
s'écrire
\begin{align}
 v_+(x,t)&=A_+\cos(2\pi f_0t-\beta x+\phi_+),\\
 v_-(x,t)&=A_-\cos(2\pi f_0t+\beta x+\phi_-).
\end{align}
Le signe devant $\beta x$ distingue les deux sens de propagation.
L'équation de ligne $\partial v/\partial x=-L'\,\partial i/\partial t$,
avec $L'$ l'inductance par unité de longueur, donne pour chaque onde
\begin{equation}
 i_+=\frac{v_+}{Z_0},\qquad i_-=-\frac{v_-}{Z_0},
 \qquad Z_0=\frac{2\pi f_0L'}{\beta}.
\end{equation}
À un même endroit et au même instant,
\begin{equation}
 v=v_++v_-,\qquad Z_0i=v_+-v_-.
\end{equation}
Le signe moins ne signifie pas que le courant réfléchi est toujours
négatif : c'est une sinusoïde. Il signifie que sa relation à la tension
est inversée lorsque le sens de propagation est inversé.

On dispose donc de la somme et de la différence de deux inconnues.
On les retrouve sans les séparer dans le temps ni en fréquence :
\begin{equation}
 \boxed{v_+=\frac{v+Z_0i}{2},\qquad v_-=\frac{v-Z_0i}{2}.}
\end{equation}
Le mécanisme d'annulation apparaît en développant :
\begin{equation}
 v-Z_0i=(v_++v_-)-(v_+-v_-)=2v_-.
\end{equation}
L'onde incidente s'annule dans cette combinaison, même si elle est présente
simultanément avec l'onde réfléchie et à la même fréquence.

\subsection{Exemple numérique à un instant donné}
Supposons, à cet instant et au même point du câble,
$v_+=1\,\mathrm V$ et $v_-=0{,}2\,\mathrm V$. Pour $Z_0=50\,\Omega$,
\begin{equation}
 v=1{,}2\,\mathrm V,\qquad
 i=\frac{1-0{,}2}{50}=0{,}016\,\mathrm A=16\,\mathrm{mA}.
\end{equation}
Les deux valeurs totales permettent de retrouver
\begin{align}
 v_+&=\frac{1{,}2+50\times0{,}016}{2}=1\,\mathrm V,\\
 v_-&=\frac{1{,}2-50\times0{,}016}{2}=0{,}2\,\mathrm V.
\end{align}
Cet exemple illustre la séparation instantanée. Le coefficient complexe
$S_{11}$ est ensuite un rapport d'amplitudes complexes, pas un rapport
instantané de sinusoïdes qui pourraient passer par zéro.

\subsection{Un pont concret : deux diviseurs de tension}
Un circuit résistif idéal permet de voir comment produire une information
de réflexion sans installer un voltmètre et un ampèremètre indépendants.
Les deux branches suivantes sont alimentées par une même tension
sinusoïdale d'amplitude complexe $E$ :
\begin{center}
\begin{tabular}{c@{\qquad\qquad}c}
\multicolumn{2}{c}{N\oe ud supérieur commun : $E$}\\
$|$ & $|$\\
$R=50\,\Omega$ & $R=50\,\Omega$\\
$|$ & $|$\\
Point A : $V_A$ & Point B : $V_B$\\
$|$ & $|$\\
$R=50\,\Omega$ & Antenne : impédance $Z$\\
$|$ & $|$\\
\multicolumn{2}{c}{Masse commune}
\end{tabular}
\end{center}
On suppose que la mesure différentielle entre A et B ne charge pas le
pont. Les lois des diviseurs de tension donnent
\begin{equation}
 V_A=\frac E2,\qquad V_B=E\frac{Z}{R+Z}.
\end{equation}
La tension de différence est donc
\begin{align}
 D=V_B-V_A
 &=E\left(\frac{Z}{R+Z}-\frac12\right)\\
 &=\frac E2\frac{Z-R}{Z+R}=\frac E2\,\Gamma.
\end{align}
Ainsi,
\begin{equation}
 \boxed{\Gamma=\frac{D}{V_A}=\frac{V_B-V_A}{V_A}.}
\end{equation}
Pour relier ce pont à la séparation tension--courant, le courant dans la
branche de droite vaut $\widetilde J=(E-V_B)/R$. Au port B,
\begin{equation}
 V^+=\frac{V_B+R\widetilde J}{2}=\frac E2=V_A,
 \qquad
 V^-=\frac{V_B-R\widetilde J}{2}=V_B-\frac E2=D.
\end{equation}
\textbf{Dans ce modèle idéal, la différence de tension du pont est
exactement l'amplitude de l'onde réfléchie au port B.}

\begin{center}
\begin{tabular}{lccc}
\toprule
Terminaison & $V_B$ & $D=V_B-V_A$ & $\Gamma$\\
\midrule
Charge $Z=R$ & $E/2$ & $0$ & $0$\\
Circuit ouvert & $E$ & $+E/2$ & $+1$\\
Court-circuit & $0$ & $-E/2$ & $-1$\\
\bottomrule
\end{tabular}
\end{center}
Le signe négatif du court-circuit correspond à un déphasage de $180^\circ$.
Le pont réalise une soustraction électrique avant la démodulation : la
voie I/Q mesure ensuite l'amplitude complexe de cette différence.

\paragraph{Limites du schéma pédagogique.}
Ce schéma n'est pas présenté comme celui du LiteVNA. À haute fréquence,
les connexions, composants, dispositifs de mesure différentielle et
couplages ont leurs propres effets. Un coupleur à lignes couplées utilise
une autre réalisation physique de la séparation directionnelle. Dans un
instrument réel, l'annulation imparfaite laisse fuir une partie de l'onde
incidente vers la voie de réflexion ; la calibration corrige cette erreur
de directivité dans les limites de stabilité et de validité du modèle.


\section{Expression idéale de \texorpdfstring{$S_{11}$}{S11} en fonction des quadratures}
Pour une mesure de réflexion à un port (ou avec les autres ports adaptés),
$S_{11}=\Gamma$. Avec les ondes normalisées
$a_1=V^+/\sqrt{Z_0}$ et $b_1=V^-/\sqrt{Z_0}$, le même rapport s'écrit
\begin{equation}
 S_{11}(f_0)=\frac{b_1}{a_1}=\frac{V^-}{V^+}.
\end{equation}
Après séparation idéale et remise à l'échelle des deux mesures,
\begin{equation}
 V^+=I_++\jj Q_+,\qquad V^-=I_-+\jj Q_-,
 \qquad
 \boxed{S_{11}(f_0)=\frac{I_-+\jj Q_-}{I_++\jj Q_+}.}
\end{equation}
En multipliant par le conjugué du dénominateur,
\begin{align}
 \Reel S_{11}&=\frac{I_-I_++Q_-Q_+}{I_+^2+Q_+^2},\\
 \operatorname{Im}S_{11}&=\frac{Q_-I_+-I_-Q_+}{I_+^2+Q_+^2},\\
 |S_{11}|&=\frac{\sqrt{I_-^2+Q_-^2}}{\sqrt{I_+^2+Q_+^2}},\\
 \arg S_{11}&=\operatorname{atan2}(Q_-,I_-)
             -\operatorname{atan2}(Q_+,I_+)\pmod{2\pi}.
\end{align}
Une rotation de phase commune aux deux mesures s'annule dans le rapport.
Une erreur différente entre les deux voies ne s'annule pas automatiquement.

\section{Du rapport brut à la mesure calibrée du banc}
\subsection{Les valeurs reçues par Python}
Dans le mode brut du pilote \texttt{plexiglas/progs/litevna.py}, les données
USB décodées fournissent, pour chaque fréquence $f_k$, deux nombres complexes :
\begin{equation}
 U_k=I_{\mathrm{ref},k}+\jj Q_{\mathrm{ref},k},\qquad
 W_k=I_{\mathrm{refl},k}+\jj Q_{\mathrm{refl},k}.
\end{equation}
Le programme calcule
\begin{equation}
 \boxed{\gamma_{m,k}=W_k/U_k.}
\end{equation}
Ces nombres sont des valeurs numériques de mesure internes, pas directement
des tensions en volts au connecteur. Leurs noms ne prouvent pas qu'ils sont
exactement égaux aux ondes idéales $V^+$ et $V^-$. Le rapport brut conserve
les erreurs de la chaîne de mesure.

Par exemple, si $U_k=1000$ et $W_k=100-100\jj$,
\begin{equation}
 \gamma_{m,k}=0{,}1-0{,}1\jj,
 \quad |\gamma_{m,k}|=0{,}1414,
 \quad \arg\gamma_{m,k}=-45^\circ.
\end{equation}
Ce résultat est encore à corriger si l'acquisition est brute.

\subsection{Calibration SOL : trois erreurs complexes par fréquence}
Le modèle utilisé dans \texttt{commun/sol.py} est
\begin{equation}
 \gamma_{m,k}=e_{00,k}+
 \frac{e_{R,k}S_{11}(f_k)}{1-e_{11,k}S_{11}(f_k)}.
\end{equation}
Les coefficients représentent la directivité résiduelle $e_{00,k}$,
le suivi en réflexion $e_{R,k}$ et la désadaptation de source $e_{11,k}$.
Dans le code, $e_R$ est enregistré sous le nom \texttt{e10e01}.
Ils sont déterminés à l'aide des trois étalons court-circuit, circuit
ouvert et charge adaptée, avec les coefficients de réflexion retenus
pour ces étalons (idéalement $-1$, $+1$ et $0$).

Posons $x_k=\gamma_{m,k}-e_{00,k}$. On résout explicitement :
\begin{align}
 x_k(1-e_{11,k}S_{11})&=e_{R,k}S_{11},\\
 x_k&=(e_{R,k}+e_{11,k}x_k)S_{11},\\
 \widehat S_{11}(f_k)
 &=\frac{\gamma_{m,k}-e_{00,k}}
 {e_{R,k}+e_{11,k}(\gamma_{m,k}-e_{00,k})}.
\end{align}
Cette calibration place la référence au plan de raccordement des étalons,
sous réserve de leur modèle et de la stabilité du montage. Le câble et
les adaptateurs situés avant ce plan sont inclus dans la correction.
L'antenne connectée ensuite n'est pas supprimée par cette calibration.

\paragraph{Autre mode du pilote.}
En mode de correction interne, le pilote attend une référence numérique
$2^{30}+\jj0$ et divise la valeur complexe reçue par $2^{30}$. Cette
normalisation numérique ne doit pas être interprétée comme une mesure
physique de tension incidente constante. Elle se distingue du rapport
brut suivi de la SOL logicielle détaillé ci-dessus.

\section{Répéter le calcul pour les \texorpdfstring{$K=201$}{K=201} fréquences}
À chaque point du balayage, $f_0$ prend successivement la valeur
$f_k=f_{\min}+k\Delta f$, $k=0,\ldots,K-1$. La chaîne conceptuelle est
\begin{equation}
 \boxed{\text{séparation RF}\ \longrightarrow\ (I,Q)\text{ de chaque mesure}
 \ \longrightarrow\ \gamma_{m,k}\ \longrightarrow\ \widehat S_{11}(f_k).}
\end{equation}
Les $K$ valeurs $\widehat S_{11}(f_k)$ sont \textbf{complexes} : on conserve
leur amplitude et leur phase. $N$, nombre éventuel d'échantillons internes
utilisés pour une démodulation, et $K$, nombre de fréquences du balayage,
sont deux quantités différentes. Le pilote ne reçoit pas les $N$
échantillons radiofréquences évoqués dans le modèle discret.

Pour obtenir ensuite un profil en retard, on peut former, par exemple,
\begin{equation}
 h(\tau)=\frac{1}{K}\sum_{k=0}^{K-1}w_k\widehat S_{11}(f_k)
 e^{+\jj2\pi f_k\tau}.
\end{equation}
Dans le traitement avec soustraction du fond, on remplace $\widehat S_{11}$
par la différence complexe entre les moyennes avec cible et sans cible.
Ici $\tau$ est un \emph{retard testé}, distinct du temps $t$ de la
sinusoïde initiale. Le profil affiché est $|h(\tau)|$ ; en réflexion, la
distance équivalente air est $d=c\tau/2$.

On ne prend donc pas le module de chaque coefficient avant de
calculer le profil de distance : cela supprimerait l'information de phase
qui porte les retards.

\section{Pourquoi un coefficient de réflexion contient-il un retard ?}
\subsection{Réflexion locale et coefficient mesuré}
Un coefficient de réflexion est un \textbf{rapport complexe}, pas seulement
une fraction d'amplitude ou de puissance. Sa phase dépend du plan où l'on
compare les ondes incidente et réfléchie. Il faut distinguer le coefficient
local $\rho$ de la cible et le coefficient observé au point de mesure.

Considérons d'abord un modèle idéal à un seul trajet, sans atténuation :
une onde se propage à la vitesse $c$ entre un plan de référence et une
surface située à la distance $d$. À la fréquence fixée $f_0$, l'onde
incidente au plan de référence a pour amplitude complexe $A$ :
\begin{equation}
 a(t)=\Reel\{A e^{\jj2\pi f_0t}\}.
\end{equation}
Après le temps de propagation $d/c$, son expression à la cible est
\begin{equation}
 a_{\mathrm{cible}}(t)
 =\Reel\{A e^{\jj2\pi f_0(t-d/c)}\},
 \qquad A_{\mathrm{cible}}=A e^{-\jj2\pi f_0d/c}.
\end{equation}
La réflexion locale multiplie cette amplitude par $\rho$ :
\begin{equation}
 B_{\mathrm{cible}}=\rho A_{\mathrm{cible}}.
\end{equation}
Par exemple, le coefficient local du champ électrique tangentiel à une
surface conductrice parfaite, à incidence normale, vaut $\rho=-1$.
Il décrit alors une inversion de phase au niveau de la surface.

Le retour vers le plan de référence ajoute un second facteur de propagation :
\begin{equation}
 B=B_{\mathrm{cible}} e^{-\jj2\pi f_0d/c}
 =\rho A e^{-\jj2\pi f_0(2d/c)}.
\end{equation}
Le rapport des deux ondes \emph{au même plan de référence} est donc
\begin{equation}
 \boxed{S_{11}(f_0)=\frac BA
 =\rho e^{-\jj2\pi f_0\tau_0},\qquad \tau_0=\frac{2d}{c}.}
\end{equation}
\textbf{Le facteur exponentiel provient du trajet aller-retour, pas de
la réflexion locale.} Un coefficient local constant n'implique donc pas
un coefficient mesuré de phase constante lorsque la fréquence varie.
Dans ce modèle, pour $\rho\ne0$ constant et une phase déroulée,
\begin{equation}
 \arg S_{11}(f)=\arg\rho-2\pi f\tau_0,\qquad
 \tau_0=-\frac{1}{2\pi}\frac{\mathrm d\arg S_{11}(f)}{\mathrm df}.
\end{equation}
La pente de phase contient l'information de retard.

\subsection{Ce que cette expression signifie avec une antenne}
Le modèle précédent est directement celui d'une ligne idéale terminée par
un réflecteur, avec $c$ remplacé par la vitesse de propagation si nécessaire.
Pour une antenne éclairant une plaque en espace libre, le couplage vers la
plaque et le retour vers le port introduisent d'autres facteurs. On écrit
la \emph{contribution de l'écho} sous la forme
\begin{equation}
 S_{\text{écho}}(f;d)=g(f,d)e^{-\jj2\pi f(2d/c)}.
\end{equation}
On note ce gain $g$ pour le distinguer de l'enveloppe temporelle
$\alpha(t)$ définie au début du document.
Le facteur complexe $g(f,d)$ comprend notamment la réflexion de la
plaque, la propagation en amplitude, la géométrie et la réponse de l'antenne.
L'approximation $g(f,d)\simeq g$ sur la bande est une hypothèse
supplémentaire. Si sa phase varie avec la fréquence, elle peut contribuer
au retard apparent et modifier la forme du pic.

Le $S_{11}$ total contient également la réflexion propre de l'antenne et
les contributions de l'environnement. Dans un modèle additif simplifié,
\begin{equation}
 S_{11}^{\mathrm{avec}}(f)=S_{\mathrm{fond}}(f)
 +S_{\text{écho}}(f;d).
\end{equation}
La soustraction complexe isole l'écho seulement si ce fond reste inchangé.
Un masquage de l'environnement par la plaque ou une dérive du montage
laisse aussi des termes résiduels. Ainsi l'expression
$g e^{-\jj2\pi f\tau_0}$ ne représente pas automatiquement tout le
$S_{11}$ mesuré.

\section{Pourquoi l'IFFT produit-elle un pic au retard de l'écho ?}
\subsection{Addition cohérente au bon retard}
On note $S_{11}^{\mathrm{corrected}}(f_k)$ le signal complexe utilisé pour
le profil : dans le traitement avec cible, il s'agit de la mesure calibrée
avec cible, après correction du ''fond'' calibré, la soustraction étant complexe.
Cette notation inclut donc ici la soustraction du fond, et pas seulement
la calibration SOL. Dans le cas idéal où ces corrections isolent un seul
écho, supposons exactement
$S_{11}^{\mathrm{corrected}}(f_k)=g e^{-\jj2\pi f_k\tau_0}$ avec $g$ constant, et prenons
une fenêtre uniforme. Pour chaque retard testé $\tau$, calculons
\begin{equation}
 h(\tau)=\frac1K\sum_{k=0}^{K-1}S_{11}^{\mathrm{corrected}}(f_k) e^{+\jj2\pi f_k\tau}
 =\frac{g}{K}\sum_{k=0}^{K-1}
 e^{\jj2\pi f_k(\tau-\tau_0)}.
\end{equation}
Si $\tau=\tau_0$, tous les facteurs exponentiels valent $1$ :
\begin{equation}
 h(\tau_0)=\frac{g}{K}(1+\cdots+1)=g.
\end{equation}
Les contributions s'ajoutent dans la même direction du plan complexe.
Pour d'autres retards, elles ont des phases différentes et se compensent
partiellement. Le module présente donc un maximum au bon retard, avec des
répétitions dues à l'échantillonnage fréquentiel.
L'IFFT recherche ainsi le retard qui remet les contributions de
 toutes les fréquences en phase.

\subsection{Forme du pic et relation exacte avec l'IFFT}
Pour $f_k=f_{\min}+k\Delta f$, posons $\delta=\tau-\tau_0$.
La somme géométrique donne
\begin{equation}
 h(\tau)=\frac{g}{K}
 e^{\jj2\pi f_{\min}\delta}
 e^{\jj\pi(K-1)\Delta f\delta}
 \frac{\sin(\pi K\Delta f\delta)}{\sin(\pi\Delta f\delta)}.
\end{equation}
Aux points où le quotient est indéterminé, sa valeur est définie par
continuité. Le maximum de module se répète tous les $1/\Delta f$.
Le pic a une largeur finie : les premiers zéros autour de $\tau_0$ sont
à $\delta=\pm1/(K\Delta f)$, soit approximativement $\pm1/B$, avec
$B=(K-1)\Delta f$. La largeur totale entre ces zéros vaut donc environ
$2/B$ ; $1/B$ est une échelle de résolution, pas cette largeur totale.

Si $z_m$ est l'IFFT de longueur $M\ge K$ de la suite $S_{11}^{\mathrm{corrected}}(f_k)$, complétée par des
zéros, selon la normalisation usuelle,
\begin{equation}
 z_m=\frac1M\sum_{k=0}^{K-1}S_{11}^{\mathrm{corrected}}(f_k) e^{\jj2\pi km/M},
 \qquad \tau_m=\frac{m}{M\Delta f},
\end{equation}
alors
\begin{equation}
 h(\tau_m)=\frac{M}{K}e^{\jj2\pi f_{\min}\tau_m}z_m,
 \qquad |h(\tau_m)|=\frac{M}{K}|z_m|.
\end{equation}
Le fait que le balayage commence à $f_{\min}>0$ introduit donc une phase
supplémentaire, mais ne déplace pas le pic du module dans ce modèle.
Le remplissage par zéros affine la grille de retard sans augmenter la
bande mesurée ni la capacité à séparer deux échos proches. Un fenêtrage
réduit les lobes secondaires, au prix d'un élargissement du pic principal.

\subsection{Conversion du retard en distance et limites physiques}
Pour un aller-retour dans l'air,
\begin{equation}
 \boxed{d=\frac{c\tau_0}{2}.}
\end{equation}
Ainsi $4\,\mathrm{ns}$ correspondent à environ $60\,\mathrm{cm}$.
Pour une bande de $1$ à $6\,\mathrm{GHz}$, $1/B=0{,}2\,\mathrm{ns}$,
soit une échelle de distance $c/(2B)\simeq3\,\mathrm{cm}$.
Cette échelle n'est pas une borne universelle de précision pour la
localisation d'un écho isolé.

Si la chaîne non corrigée ajoute un retard $\tau_a$, l'observation devient
$g e^{-\jj2\pi f(2d/c+\tau_a)}$ et le pic se place à
\begin{equation}
 d_{\mathrm{apparent}}=d+\frac{c\tau_a}{2}.
\end{equation}
Une SOL au bout du câble corrige la chaîne située avant son plan de
référence ; elle ne retire pas la réponse de l'antenne connectée ensuite.
Une réponse dispersive de l'antenne ou plusieurs trajets peuvent aussi
déformer le profil : leur effet ne se résume pas nécessairement à un
décalage constant.


\section*{Conclusion}
La démodulation transforme une tension sinusoïdale en un nombre complexe
$I+\jj Q$. Elle ne sépare pas les contributions de même fréquence.
Pour obtenir $S_{11}$, il faut soit connaître l'onde incidente complexe,
soit disposer d'une information directionnelle supplémentaire. Le rapport
réfléchi/incident définit alors $S_{11}$ ; dans le banc réel, la calibration
corrige le rapport numérique mesuré.

\begin{thebibliography}{9}
\bibitem{keysight} Keysight, \emph{Network Analyzer Basics},
\href{https://www.keysight.com/us/en/assets/3119-1082/application-notes/5965-7708.pdf}{document 5965-7708}.
Principe de séparation des signaux incidents et réfléchis.
\bibitem{litevna} \emph{LiteVNA User Guide},
\href{https://nanovna.com/wp-content/uploads/2021/11/LiteVNA_User-Guide.pdf}{manuel du LiteVNA}.
Mélangeur unique et commutation RF ; ce document ne constitue pas ici une
vérification du schéma détaillé du circuit directionnel.
\end{thebibliography}
\end{document}
